\subsection{導入：なぜセキュリティを開発で扱うのか}
コンピュータシステムは、誤用、悪意ある活動、金銭目的の攻撃、情報漏えい、サービス停止の標的になる。したがって、開発者は通常の正しさや信頼性だけでなく、攻撃されても守るべきものを守れるかを考えなければならない。

安全なソフトウェア開発とは、確立された規則や推奨実践に従い、セキュリティを作り込むことである。安全な保守とは、保守変更で新しいセキュリティ問題を入れないこと、発見された脆弱性をライフサイクル中に扱えるようにすることである。

第13章は、セキュリティを「最後に診断する項目」ではなく、「要求、設計、構築、テスト、保守、組織運営に埋め込む活動」として説明している。

\begin{notebox}{重要}セキュリティ上の失敗は、機能が動かない失敗とは違う。普段は問題なく動くソフトウェアでも、悪意ある入力、権限の濫用、第三者部品の脆弱性、設定不備によって破られることがある。\end{notebox}
